\documentclass[aps,preprint]{revtex4}%
\usepackage{amssymb}
\usepackage{amsfonts}
\usepackage{amsmath}
\usepackage{graphicx}%
\setcounter{MaxMatrixCols}{30}
%TCIDATA{OutputFilter=latex2.dll}
%TCIDATA{Version=5.50.0.2960}
%TCIDATA{CSTFile=revtex4.cst}
%TCIDATA{Created=Friday, July 19, 2019 09:14:18}
%TCIDATA{LastRevised=Wednesday, January 15, 2020 11:14:57}
%TCIDATA{<META NAME="GraphicsSave" CONTENT="32">}
%TCIDATA{<META NAME="SaveForMode" CONTENT="1">}
%TCIDATA{BibliographyScheme=Manual}
%TCIDATA{<META NAME="DocumentShell" CONTENT="Articles\SW\REVTeX 4">}
%TCIDATA{Language=American English}
%BeginMSIPreambleData
\providecommand{\U}[1]{\protect\rule{.1in}{.1in}}
%EndMSIPreambleData
\newtheorem{theorem}{Theorem}
\newtheorem{acknowledgement}[theorem]{Acknowledgement}
\newtheorem{algorithm}[theorem]{Algorithm}
\newtheorem{axiom}[theorem]{Axiom}
\newtheorem{claim}[theorem]{Claim}
\newtheorem{conclusion}[theorem]{Conclusion}
\newtheorem{condition}[theorem]{Condition}
\newtheorem{conjecture}[theorem]{Conjecture}
\newtheorem{corollary}[theorem]{Corollary}
\newtheorem{criterion}[theorem]{Criterion}
\newtheorem{definition}[theorem]{Definition}
\newtheorem{example}[theorem]{Example}
\newtheorem{exercise}[theorem]{Exercise}
\newtheorem{lemma}[theorem]{Lemma}
\newtheorem{notation}[theorem]{Notation}
\newtheorem{problem}[theorem]{Problem}
\newtheorem{proposition}[theorem]{Proposition}
\newtheorem{remark}[theorem]{Remark}
\newtheorem{solution}[theorem]{Solution}
\newtheorem{summary}[theorem]{Summary}
\newenvironment{proof}[1][Proof]{\noindent\textbf{#1.} }{\ \rule{0.5em}{0.5em}}
\begin{document}
\preprint{HEP/123-qed}
\title[Short title for running header]{REV\TeX4}
\author{First Author}
\affiliation{My Institution}
\author{Second Author}
\affiliation{My Institution}
\author{Third Author}
\affiliation{Other Institution}
\keywords{one two three}
\pacs{PACS number}

\begin{abstract}
Shell document for REV\TeX{} 4.

\end{abstract}
\volumeyear{year}
\volumenumber{number}
\issuenumber{number}
\eid{identifier}
\date[Date text]{date}
\received[Received text]{date}

\revised[Revised text]{date}

\accepted[Accepted text]{date}

\published[Published text]{date}

\startpage{101}
\endpage{102}
\maketitle
\tableofcontents


\section{Lie Algebras}

A \emph{Lie algebra} is a vector space $\mathcal{G}$ together with a
\emph{non-associative} antisymmetric product called the Lie bracket, which is
an alternating bilinear map
\begin{gather}
\left[  ,\right]  :\mathcal{G}\times\mathcal{G}\rightarrow\mathcal{G}\\
\left[  ,\right]  :\left(  x,y\right)  \rightarrow\left[  x,y\right]
\end{gather}
satisfying the \emph{Jacobi Identity}%
\begin{equation}
\left[  x,\left[  y,z\right]  \right]  =\left[  \left[  x,y\right]  ,z\right]
+\left[  y,\left[  x,z\right]  \right]  .
\end{equation}


In particular
\begin{align}
\left[  ,\right]   &  :\left(  0,y\right)  \rightarrow\left[  0,y\right]  =0\\
\left[  ,\right]   &  :\left(  x,0\right)  \rightarrow\left[  x,0\right]  =0\\
\left[  ,\right]   &  :\left(  x,x\right)  \rightarrow\left[  x,x\right]  =0.
\end{align}


\begin{remark}
Lie algebras have no unit element as
\begin{equation}
\left[  I,x\right]  =x
\end{equation}
contradicts the Jacobi Identity unless $x=0.$
\end{remark}

\begin{condition}
In infinite dimensions we assume that $\mathcal{G}$ is a Banach space with
norm $\left\Vert {}\right\Vert _{\mathcal{G}}$ and the Lie bracket is
continous wrt the norm.
\end{condition}

As $\mathcal{G}$ is a vector space one can consider the \emph{associative
algebra} of linear \emph{endomorphisms} $\mathcal{L}\left(  \mathcal{G}%
\right)  $ acting in $\mathcal{G}$ and the \emph{general linear group}
$GL\left(  \mathcal{G}\right)  $ of invertible elements of $\mathcal{L}\left(
\mathcal{G}\right)  .$ The group $GL\left(  \mathcal{G}\right)  $ is a Lie
group and its Lie algebra, $gl\left(  GL\left(  \mathcal{G}\right)  \right)
,$ is $\mathcal{L}\left(  \mathcal{G}\right)  $ i.e. $\mathcal{L}\left(
\mathcal{G}\right)  $ is closed under the multiplication
\begin{gather}
\left[  ,\right]  :\mathcal{L}\left(  \mathcal{G}\right)  \times
\mathcal{L}\left(  \mathcal{G}\right)  \rightarrow\mathcal{L}\left(
\mathcal{G}\right) \\
\left[  ,\right]  :\left(  A,B\right)  \rightarrow\left[  A,B\right] \\
\forall A,B\in\mathcal{L}\left(  \mathcal{G}\right)  .
\end{gather}


One can also consider the \emph{exterior algebra} $%
%TCIMACRO{\tbigwedge }%
%BeginExpansion
{\textstyle\bigwedge}
%EndExpansion
\left(  \mathcal{G}\right)  $ defined in the standard manner by
\begin{equation}%
%TCIMACRO{\tbigwedge }%
%BeginExpansion
{\textstyle\bigwedge}
%EndExpansion
\left(  \mathcal{G}\right)  =%
%TCIMACRO{\tbigoplus \limits_{0\leq n\leq\infty}}%
%BeginExpansion
{\textstyle\bigoplus\limits_{0\leq n\leq\infty}}
%EndExpansion%
%TCIMACRO{\tbigwedge \nolimits^{n}}%
%BeginExpansion
{\textstyle\bigwedge\nolimits^{n}}
%EndExpansion
\mathcal{G}.
\end{equation}


\subsubsection{Enveloping Algebra}

The \emph{associative enveloping algebra} called the \emph{multiplication
algebra, }$\mathcal{M}\left(  \mathcal{G}\right)  ,$ of $\mathcal{G}$ is the
associative algebra generated by the \emph{left and right }linear
multiplication maps.

\begin{definition}
The \emph{Centroid of} $\mathcal{G}$ is the \emph{Centraliser} of the
\emph{Enveloping Algebra} in the \emph{endomorphism algebra} $\mathcal{L}%
\left(  \mathcal{G}\right)  $.
\end{definition}

\begin{definition}
An algebra is \emph{Central} if its \emph{centroid }consists of the
$\mathbb{C}$-multiples of the identity of $\mathcal{L}\left(  \mathcal{G}%
\right)  .$
\end{definition}

\subparagraph{Left Multiplication SubAlgebra}

The left multiplication operator $L_{A}\in\mathcal{L}\left(  \mathcal{G}%
\right)  $ is defined by
\begin{equation}
L_{A}\left(  B\right)  \equiv L\left(  A\right)  \left(  B\right)  =\left[
A,B\right]  \,\forall B\in\mathcal{G},
\end{equation}
leading to
\begin{equation}
L_{A}^{2}\left(  B\right)  =\left[  A,\left[  A,B\right]  \right]  \,\forall
B\in\mathcal{G}.
\end{equation}
Left multiplication forms a sub associative algebra, $L_{m}\left(
\mathcal{G}\right)  ,$ of $\mathcal{L}\left(  \mathcal{G}\right)  $ and a sub
Lie algebra, $L\left(  \mathcal{G}\right)  ,$ of $\mathcal{L}\left(
\mathcal{G}\right)  $ as
\begin{gather}
\left(  L_{A}\circ L_{B}\right)  \left(  C\right)  =\left[  A,\left[
B,C\right]  \right]  \,=A\left(  BC-CB\right)  -\left(  BC-CB\right)
A=ABC-ACB-BCA+CBA\\
\left(  L_{B}\circ L_{A}\right)  \left(  C\right)  =\left[  B,\left[
A,C\right]  \right]  =B\left(  AC-CA\right)  -\left(  AC-CA\right)
B=BAC-BCA-ACB+CAB.
\end{gather}
giving%
\begin{align}
\left(  \left(  L_{A}\circ L_{B}\right)  -\left(  L_{B}\circ L_{A}\right)
\right)  \left(  C\right)   &  =ABC-ACB-BCA+CBA-BAC+BCA+ACB-CAB\\
&  =ABC+CBA-BAC-CAB=\left[  A,B\right]  C-C\left[  A,B\right] \\
&  =\left[  \left[  A,B\right]  ,C\right]  .
\end{align}
Thus one obtains
\begin{equation}
\left[  L_{A},L_{B}\right]  \left(  C\right)  =\left[  L_{\left[  A,B\right]
},C\right]  \quad\forall C\in\mathcal{G},
\end{equation}
hence
\begin{equation}
L:\mathcal{G}\rightarrow L\left(  \mathcal{G}\right)  \subset\mathcal{L}%
\left(  \mathcal{G}\right)
\end{equation}
is a Lie algebra representation.

\begin{remark}
$L_{A}$ is a \emph{Derivation}, i.e. $L_{A}\in\operatorname{Der}\left\{
\mathcal{G}\right\}  $ acting in $\mathcal{G}$ i.e
\begin{equation}
L_{A}\left(  \left[  B,C\right]  \right)  =\left[  L_{A}\left(  B\right)
,C\right]  +\left[  B,L_{A}\left(  C\right)  \right]  .
\end{equation}

\end{remark}

\begin{remark}
The associative subalgebra $L_{m}\left(  \mathcal{G}\right)  \subset
\mathcal{L}\left(  \mathcal{G}\right)  $ is an \emph{enveloping algebra }of
$\mathcal{G}.$
\end{remark}

\subparagraph{Right Multiplication SubAlgebra}

In an analogous manner to left multplication one defines the right
multiplication operators $R_{B}\in\mathcal{L}\left(  \mathcal{G}\right)  $ by
\begin{equation}
R_{B}\left(  A\right)  \equiv R\left(  B\right)  \left(  A\right)  =\left[
A,B\right]  \,\forall A\in\mathcal{G},
\end{equation}
leading to
\begin{equation}
R_{A}^{2}\left(  B\right)  =\left[  \left[  A,B\right]  ,B\right]  \,\forall
A\in\mathcal{G}.
\end{equation}
Right multiplication forms a sub algebra, $R_{m}\left(  \mathcal{G}\right)  ,$
of $\mathcal{L}\left(  \mathcal{G}\right)  ,$ however $R\left(  \mathcal{G}%
\right)  $ is \emph{not a represntation of the Lie algebra }$\mathcal{G}$ as
\begin{gather}
\left(  R_{A}\circ R_{B}\right)  \left(  C\right)  =\left[  \left[
C,B\right]  ,A\right]  \,=\left(  CB-BC\right)  A-A\left(  CB-BC\right)
=CBA-BCA-ACB+ABC\\
\left(  R_{B}\circ R_{A}\right)  \left(  C\right)  =\left[  \left[
C,A\right]  ,B\right]  =\left(  CA-AC\right)  B-B\left(  CA-AC\right)
=CAB-ACB-BCA+BAC
\end{gather}
giving
\begin{align}
\left[  R_{A},R_{B}\right]   &  =\left(  R_{A}\circ R_{B}-R_{B}\circ
R_{A}\right)  \left(  C\right)  \\
&  =CBA-BCA-ACB+ABC-CAB+ACB+BCA-BAC\\
&  =CBA+ABC-CAB-BAC\\
&  =C\left[  B,A\right]  -\left[  B,A\right]  C=\left[  C,\left[  B,A\right]
\right]  =R_{\left[  B,A\right]  }.
\end{align}
Thus
\begin{equation}
R:\mathcal{G}\rightarrow\mathcal{L}\left(  \mathcal{G}\right)
\end{equation}
is \emph{NOT} a Lie algebra representation as
\begin{equation}
R_{\left[  A,B\right]  }\neq R_{A}\circ R_{B}%
\end{equation}
i.e. the order of the product is reversed.

\begin{remark}
The associative subalgebra $R_{m}\left(  \mathcal{G}\right)  \subset
\mathcal{L}\left(  \mathcal{G}\right)  $ is an \emph{enveloping algebra} of
$\mathcal{G}.$
\end{remark}

\begin{remark}
Right and left multiplication commute i.e
\begin{equation}
R\left(  A\right)  L\left(  B\right)  =L\left(  B\right)  R\left(  A\right)
\quad\forall A,B\in\mathcal{G}%
\end{equation}
Thus in particular the multiplicative algebra $\mathcal{M}\left(
\mathcal{G}\right)  $ is a \emph{Flexible algebra.}
\end{remark}

\begin{remark}
The associative subalgebra $\mathcal{M}\left(  \mathcal{G}\right)  $ generated
by $\left\{  L\left(  \mathcal{G}\right)  ,R\left(  \mathcal{G}\right)
\right\}  \subset\mathcal{L}\left(  \mathcal{G}\right)  $ is an
\emph{enveloping algebra} of $\operatorname{adj}\left\{  \mathcal{G}\right\}
\subset\mathcal{B}\left(  \mathcal{L}\left(  \mathcal{G}\right)  \right)  .$
\end{remark}

\begin{definition}
A \emph{Lie subalgebra} of $\mathcal{G}$ is a subspace $\mathcal{H}$ such
that
\begin{equation}
\left[  \mathcal{H},\mathcal{H}\right]  \subseteq\mathcal{H}%
\end{equation}

\end{definition}

\begin{definition}
A \emph{Lie ideal}, $\mathcal{H},$ of $\mathcal{G}$ is a subspace
$\mathcal{H}$ such that
\begin{equation}
\left[  \mathcal{G},\mathcal{H}\right]  \subseteq\mathcal{H}.
\end{equation}

\end{definition}

Let $\mathcal{H}$ and $\mathcal{K}$ be subspaces of $\mathcal{G}$ then we can
form the subspace $\ $
\begin{equation}
\left[  \mathcal{H},\mathcal{K}\right]  =\operatorname*{linspan}\left\{
\left.  \left[  h,k\right]  \right\vert h\in\mathcal{H},k\in\mathcal{K}%
\right\}  .
\end{equation}


\begin{definition}
The \emph{Center}, $\mathcal{Z}\left(  \mathcal{G}\right)  ,$ of the Lie
algebra
\begin{equation}
\mathcal{Z}\left(  \mathcal{G}\right)  =\left\{  \left.  x\right\vert \left[
x,y\right]  =0\,\forall y\in\mathcal{G}\right\}
\end{equation}

\end{definition}

\begin{definition}
The \emph{Derived algebra} $\left[  \mathcal{G},\mathcal{G}\right]  $
\begin{equation}
\left[  \mathcal{G},\mathcal{G}\right]  =\operatorname*{linspan}\left\{
\left.  \left[  h,k\right]  \right\vert h\in\mathcal{G},k\in\mathcal{G}%
\right\}
\end{equation}
is an ideal of $\mathcal{G}.$
\end{definition}

\begin{definition}
If $\mathcal{G}$ contains only trivial ideals ($\mathcal{G}$ and $0$) then it
is \emph{Simple.}
\end{definition}

\begin{lemma}
\begin{enumerate}
\item If $\mathcal{H}$ and $\mathcal{K}$ are subalgebras then $\mathcal{H}%
\cap\mathcal{K}$ is a subalgebra

\item If $\mathcal{H}$ and $\mathcal{K}$ are ideals then $\mathcal{H}%
\cap\mathcal{K}$ is an ideal

\item If $\mathcal{H}$ is an ideal and $\mathcal{K}$ is a subalgebras then
$\mathcal{H}+\mathcal{K}$ is a subalgebra$.$
\end{enumerate}
\end{lemma}

The \emph{Centralizer}, $\mathcal{C}_{L}\left(  \mathcal{H}\right)  ,$ of
$\mathcal{H}$ in $\mathcal{G}$ is defined by
\begin{equation}
\mathcal{C}_{L}\left(  \mathcal{H}\right)  =\left\{  \left.  k\in
\mathcal{G}\right\vert \left[  \mathcal{H},k\right]  =0\right\}
\end{equation}
and the \emph{Normalizer, }$\mathcal{N}_{L}\left(  \mathcal{H}\right)
,$\emph{ }%
\begin{equation}
\mathcal{N}_{L}\left(  \mathcal{H}\right)  =\left\{  \left.  k\in
\mathcal{G}\right\vert \left[  \mathcal{H},k\right]  \subseteq\mathcal{H}%
\right\}  .
\end{equation}


\paragraph{Exponential function on $gl\left(  GL\left(  \mathcal{G}\right)
\right)  $}

The \emph{exponential function} maps the left multiplication representation
$L\left(  \mathcal{G}\right)  $ of $\mathcal{G}$ \underline{\emph{onto}} the
Lie group $GL\left(  \mathcal{G}\right)  \subset\mathcal{M}\left(
\mathcal{G}\right)  \subset\mathcal{L}\left(  \mathcal{G}\right)  $ i.e. it is
surjective
\begin{align}
\exp &  :L\left(  \mathcal{G}\right)  \rightarrow GL\left(  \mathcal{G}%
\right)  \\
e^{A} &  \in GL\left(  \mathcal{G}\right)  \,\forall A\in L\left(
\mathcal{G}\right)
\end{align}
furthermore one can show that for this general linear group as the mapping is
surjective
\begin{align}
A  & =\log\left\{  X\right\}  =\log\left\{  \exp\left\{  A\right\}  \right\}
\in L\left(  \mathcal{G}\right)  \\
X  & =\exp\left\{  \log\left\{  X\right\}  \right\}  \in GL\left(
\mathcal{G}\right)  .
\end{align}


\begin{proposition}
To every Lie algebra $\mathcal{G}$ there is a unique simply connected Lie
group $G$ with Lie algebra $\mathcal{G}$ hence the universal covering group of
a connected Lie group $H$ is the \emph{(unique)} simply connected Lie group
$G$ having the same Lie algebra as $H$.
\end{proposition}

\begin{remark}
In general $G\left(  \mathcal{H}\right)  \supset\exp\left\{  L\left(
\mathcal{H}\right)  \right\}  $ and $\operatorname{Ker}\left\{  \left.
\exp\right\vert _{\mathcal{H}}\right\}  \neq0$ i.e $\exp$ \emph{might not be
surjective nor injective}, where $\mathcal{H}$ is a sub Lie algebra of
$\mathcal{G}$ and $G\left(  \mathcal{H}\right)  $ is a group that
\emph{covers} $\exp\left\{  L\left(  \mathcal{H}\right)  \right\}  .$
\end{remark}

The $\log$ series
\begin{align}
\log\left\{  B\right\}    & =\sum_{1\leq k\leq\infty}\frac{\left(  -1\right)
^{k+1}\left(  B-I\right)  ^{k}}{k}\\
\log\left\{  I+A\right\}    & =\sum_{1\leq k\leq\infty}\frac{\left(
-1\right)  ^{k+1}A^{k}}{k}%
\end{align}
converges when $\left\Vert \left(  B-I\right)  \right\Vert _{\mathcal{L}%
\left(  \mathcal{G}\right)  }<1\equiv\left\Vert A\right\Vert _{\mathcal{L}%
\left(  \mathcal{G}\right)  }<1$ and defines the inverse of the exponential
\begin{equation}
e^{\log\left\{  B\right\}  }=B.
\end{equation}


\begin{remark}
\emph{A complex matrix has a logarithm if and only if it is invertible,
however the logarithm is not unique}. If a matrix \emph{has no negative real}
eigenvalues, then there is a \emph{unique logarithm} that has eigenvalues all
lying in the strip
\begin{equation}
\{z\in\mathbb{C}|-\U{3c0} <Imz<\U{3c0} \}.
\end{equation}
This logarithm is known as the \emph{principal logarithm. }
\end{remark}

\begin{remark}
A real matrix has a real logarithm if and only if it is invertible and each
Jordan block belonging to a negative eigenvalue occurs an even number of
times. If an invertible real matrix does not satisfy this condition with the
Jordan blocks, then it has only non-real logarithms. This can already be seen
in the scalar case: no branch of the logarithm can be real at $-1$.
\end{remark}

\subsubsection{Determinant}

\begin{definition}
The \emph{determinant of a trace class} operator $A\in\mathcal{L}%
_{\operatorname{Tr}}\left(  \mathcal{G}\right)  $ can be defined by
\end{definition}

%

\begin{equation}
\det\left\{  \exp\left\{  A\right\}  \right\}  =\exp\left\{  \operatorname{Tr}%
\left\{  A\right\}  \right\}  \neq0.
\end{equation}
In particular this defines the determinat of any $\exp\left\{  A\right\}  \in$
$L_{\operatorname{Tr}}\left(  \mathcal{G}\right)  \subset\mathcal{L}%
_{\operatorname{Tr}}\left(  \mathcal{G}\right)  .$

\subparagraph{Fredholm Determinants}

This definition can be extended to the set of \emph{Fredholm} operators
$F_{\mathcal{L}\left(  \mathcal{G}\right)  }$ defined by
\begin{equation}
F_{\mathcal{L}\left(  \mathcal{G}\right)  }=\left\{  \left.  I_{\mathcal{L}%
\left(  \mathcal{G}\right)  }+A\right\vert A\in\mathcal{L}_{\operatorname{Tr}%
}\left(  \mathcal{G}\right)  \right\}  .
\end{equation}
The set $F_{\mathcal{L}\left(  \mathcal{G}\right)  }$ is \emph{a group} as the
inverse is given by%
\begin{equation}
\left(  I_{\mathcal{L}\left(  \mathcal{G}\right)  }+A\right)  ^{-1}%
=I_{\mathcal{L}\left(  \mathcal{G}\right)  }-A\left(  I_{\mathcal{L}\left(
\mathcal{G}\right)  }+A\right)  ^{-1}\in F_{\mathcal{L}\left(  \mathcal{G}%
\right)  }.
\end{equation}
and
\begin{equation}
A\left(  I_{\mathcal{L}\left(  \mathcal{G}\right)  }+A\right)  ^{-1}%
\in\mathcal{L}_{\operatorname{Tr}}\left(  \mathcal{G}\right)  .
\end{equation}


\begin{remark}
In general $\left(  I_{\mathcal{L}\left(  \mathcal{G}\right)  }+A\right)
\notin\mathcal{L}_{\operatorname{Tr}}\left(  \mathcal{G}\right)  .$
\end{remark}

\begin{definition}
If $\mathcal{G}$ has a basis then one can also define the \emph{Fredholm
Determinant} as
\begin{equation}
\det\left\{  \left(  I_{\mathcal{L}\left(  \mathcal{G}\right)  }+A\right)
\right\}  =%
%TCIMACRO{\tsum \limits_{0\leq n\leq\infty}}%
%BeginExpansion
{\textstyle\sum\limits_{0\leq n\leq\infty}}
%EndExpansion
\operatorname{Tr}\left\{  \boxtimes^{n}A\right\}
\end{equation}
when $A\in\mathcal{L}_{\operatorname{Tr}}\left(  \mathcal{G}\right)  $ and we
have extended the linear functional $\operatorname{Tr}_{\mathcal{L}\left(
\mathcal{G}\right)  }$ to the linear functional $\operatorname{Tr}%
_{\mathcal{L}\left(  \wedge\left(  \mathcal{G}\right)  \right)  }$ acting in
$\mathcal{L}_{\operatorname{Tr}}\left(  \bigwedge\left(  \mathcal{G}\right)
\right)  .$
\end{definition}

This definition leads to
\begin{equation}
\det\left\{  \left(  I_{\mathcal{L}\left(  \mathcal{G}\right)  }+zA\right)
\right\}  =%
%TCIMACRO{\tsum \limits_{0\leq n\leq\infty}}%
%BeginExpansion
{\textstyle\sum\limits_{0\leq n\leq\infty}}
%EndExpansion
z^{n}\operatorname{Tr}\left\{  \boxtimes^{n}A\right\}  ;\quad z\in\mathbb{C}.
\end{equation}


\subsubsection{Topology, Homotopy, Lie Groups and Algebras.}

\begin{definition}
In \emph{topology} two \emph{continuous functions} from one topological space
to another are called \emph{homotopic} if one can be "continuously deformed"
into the other, such a deformation are called a \emph{homotopy} between the
two functions.
\end{definition}

\begin{definition}
A topological space $X$ is said to be \emph{path-connected or arc-wise
connected} if given any two points $\left\{  a,b\right\}  $in the topological
space, there is a \emph{continous path} $\gamma\in C\left(  \left[
0,1\right]  ,X\right)  $
\begin{align}
\gamma &  :\left[  0,1\right]  \rightarrow X\\
\gamma\left(  0\right)   &  =a\in X,\quad\gamma\left(  1\right)  =b\in X.
\end{align}

\end{definition}

\begin{definition}
A topological space is said to be \emph{simply connected} if it satisfies the
following equivalent conditions:
\end{definition}

\begin{enumerate}
\item It is \emph{path-connected}, and \emph{any loop (path with same initial
and endpoint) at any point} is \emph{homotopic to the constant loop at that
point}

\item It is \emph{path-connected}, and its \emph{fundamental group is
trivial.}
\end{enumerate}

\begin{definition}
A topological space $X$ is called \emph{simply connected} if it is
\emph{path-connected} and any \emph{loop} $f$ in $X$ defined by
\begin{equation}
f:S^{1}\rightarrow X
\end{equation}
can be contracted to a point.
\end{definition}

\begin{definition}
A topological group is termed a \emph{simply connected group} if its
underlying topological space is a simply connected space.
\end{definition}

\begin{definition}
The \emph{fundamental group} $\Pi_{1}\left(  X\right)  $ of an
\emph{arcwise-connected} topological space $X$ is the \emph{group formed} by
the \emph{sets of equivalence classes} of the \emph{set of all loops}, i.e.,
\emph{paths} with the same initial and final points at a given basepoint $p$,
under the \emph{equivalence relation of homotopy}. The \emph{identity element}
of this group is the set of all \emph{paths homotopic} to the \emph{degenerate
path consisting of the point }$p$. The fundamental groups of homeomorphic
spaces are isomorphic. In fact, the fundamental group only depends on the
\emph{homotopy type} of $X$.
\end{definition}

\begin{definition}
A \emph{Connected Lie Group} is a \emph{path-connected topological space.}
\end{definition}

\begin{definition}
A \emph{Simply Connected} Lie Group is a \emph{simply connected}
\emph{topological space}
\end{definition}

\begin{definition}
A \emph{Covering Group} of a topological group $H$ is a \emph{covering space}
$G$ of $H$ such that $G$ is a topological group and the covering map%
\begin{equation}
\rho:G\rightarrow H
\end{equation}
is a \emph{continuous group homomorphism}. The map $\rho$ is called the
\emph{covering homomorphism}.
\end{definition}

\begin{remark}
Group homomorphisms are in $1-1$ correspondence with \emph{normal} subgroups. 
\end{remark}

\begin{definition}
A covering space is a \emph{Universal Covering Space} if it is simply connected.
\end{definition}

\begin{definition}
The \emph{index} of a group $H$ in $G$ is defined as the number of cosets of
$H$ in $G$. (It is always the case that the number of left cosets of $H$ in
$G$ is equal to the number of right cosets.) The index is denoted by
$\left\vert G:H\right\vert .$
\end{definition}

Let $G$ be a covering group of $H$. The kernel $K$ of the covering
homomorphism $\rho$ is just the \emph{fiber over the identity} in $H,$ i.e.
$\rho^{-1}\left(  id_{H}\right)  ,$ and is a \emph{discrete normal subgroup}
of $G$. The kernel $K$ is closed in $G$ if and only if $G$ is Hausdorff (and
if and only if $H$ is Hausdorff). Going in the other direction, if $G$ is any
topological group and $K$ is a \emph{discrete normal subgroup} of $G$ then the
quotient map
\begin{equation}
\rho:G\rightarrow H\equiv G/K
\end{equation}
is a covering homomorphism. Given any \emph{discrete} \emph{normal subgroup}
of a Lie group the quotient group is a Lie group and the \emph{quotient map}
is a \emph{covering} homomorphism.

\begin{remark}
The index is the order of the cover.
\end{remark}

\begin{definition}
Two \emph{Lie groups} are \emph{locally isomorphic} if and only if their Lie
algebras are isomorphic.
\end{definition}

\begin{lemma}
A homomorphism
\begin{equation}
\rho:G\rightarrow H
\end{equation}
of Lie groups is a covering homomorphism if and only if the induced map on Lie
algebras%
\begin{equation}
\rho_{\ast}:\mathcal{G}\rightarrow\mathcal{H}%
\end{equation}
is an isomorphism.
\end{lemma}

\emph{We have the important theorem and corollary:}

\begin{theorem}
For every Lie algebra $\mathcal{G}$ there is a unique simply connected Lie
group $G$ with Lie algebra $\mathcal{G}$.
\end{theorem}

\begin{corollary}
The \emph{universal covering group} of a \emph{connected Lie group} $H$ is the
(unique) \emph{simply connected} Lie group $G$ having the same Lie algebra as
$H$.
\end{corollary}

\begin{remark}
The Lie algebra \emph{only sees the} \emph{connected subgroup that contains
the identity}. The map, $\mathcal{X},$ from the Lie group $G$ to the Lie
algebra $\mathcal{G}$
\begin{equation}
\mathcal{X}:G\rightarrow\mathcal{G}%
\end{equation}
is a lineariztion of the group around the identity and is in fact a map from
the group into the \emph{tangent space} at the identity. As smooth curves
passing through the identity only exist with the largest connected subgroup
containing the identity $\equiv$ connected component of the group containg the
identity, $\mathcal{X}$ is essentailly a map from this subgroup to the tangent
space at the identity $\mathcal{T}_{I_{G}}\equiv$ Lie algebra $\mathcal{G}.$
\emph{Thus the Lie algebra is the Lie algebra of this connected subgroup and
the Lie algebra of the largest connected subgroup containing the identity is
the same as the Lie algebra of the containing group.}
\end{remark}

\begin{example}
The general linear group over the field of complex numbers, $GL(n,\mathbb{C}%
)$, is a complex Lie group of complex dimension $n^{2}$. As a real Lie group
(through \emph{realification}) it has dimension $4n^{2}$. The set of all real
matrices forms a real Lie subgroup. These correspond to the inclusions%
\begin{equation}
GL(n,\mathbb{R})\subset GL(n,\mathbb{C})\subset GL(2n,\mathbb{R}),
\end{equation}
which have real dimensions $n^{2}$, $2n^{2}$, and $4n^{2}$ $=(2n)^{2}$.
Complex $n$-dimensional matrices can be characterized as real $2n$-dimensional
matrices that preserve a linear complex structure, i.e. that \emph{commute
with a matrix} $J$ such that $J^{2}$ $=-I$, where $J$ corresponds to
multiplying by the imaginary unit $i$.

The Lie algebra corresponding to $GL(n,\mathbb{C})$ consists of all $n\times
n$ complex matrices with the commutator serving as the Lie bracket. 
\end{example}

\begin{itemize}
\item Unlike the real case, $GL(n,\mathbb{C})$ \emph{is connected}. This
follows, in part, since the multiplicative group of complex numbers
$C^{\times}$ is connected.

\item The group manifold $GL(n,\mathbb{C})$ is \emph{not compact}; rather its
\emph{maximal compact subgroup} is the unitary group $U(n)$. 

\item As for $U(n)$, the group manifold $GL(n,\mathbb{C})$ is \emph{not simply
connected} but has a \emph{fundamental group} isomorphic to $\mathbb{Z}$.
\end{itemize}

\begin{remark}
The eponential map
\begin{equation}
\exp:gl(n,\mathbb{C})=M\left(  n,\mathbb{C}\right)  \rightarrow
GL(n,\mathbb{C})
\end{equation}
is surjective, however the exponential map
\begin{equation}
\exp:gl(n,\mathbb{C})=M\left(  n,\mathbb{C}\right)  \rightarrow
GL(n,\mathbb{C})
\end{equation}
\emph{is not injective}.
\end{remark}

\paragraph{Lie Algebra Adjoint Representation}

As $\mathcal{G}$ is a vector space this product can be used to form the
\emph{adjoint}$\equiv$\emph{regular} representation of $\mathcal{G}$
\begin{align}
ad  &  :\mathcal{G}\rightarrow ad\left(  \mathcal{G}\right)  \subset
\mathcal{M\left(  \mathcal{G}\right)  }\subset\mathcal{L}\left(
\mathcal{G}\right) \\
ad\left(  x\right)  \left(  y\right)   &  =\left[  x,y\right]  =z.
\end{align}
In particular this gives
\begin{equation}
ad\left(  x\right)  \left(  x\right)  =ad\left(  0\right)  \left(  x\right)
=ad\left(  x\right)  \left(  0\right)  =0\forall\,x\in\mathcal{G}.
\end{equation}
The operator $ad\left(  x\right)  \in\operatorname{Der}\left\{  \mathcal{G}%
\right\}  \subset$ $\mathcal{L}\left(  \mathcal{G}\right)  $ and is an
\emph{inner derivation}. It can be shown that if $\mathcal{G}$ is
$\emph{semisimple}$ then all derivations are inner.

If the center, $\mathcal{Z}\left(  \mathcal{G}\right)  ,$ is \emph{non
trivial} then
\begin{equation}
\ker\left\{  ad\right\}  \neq\left\{  0\right\}
\end{equation}
and $ad$ is not injective.

One can see that the operator $ad\left(  x\right)  $ is a linear operator by
\begin{align}
ad\left(  x\right)  \left(  \alpha y+\beta z\right)   &  =\left[  x,\alpha
y+\beta z\right] \\
&  =\alpha\left[  x,y\right]  +\beta\left[  x,z\right] \\
&  =\alpha\left(  ad\left(  x\right)  \left(  y\right)  \right)  +\beta\left(
ad\left(  x\right)  \left(  z\right)  \right)
\end{align}
and the set of such linear operators is closed under linear combinations i.e.
\[
\left(  \alpha ad\left(  x\right)  +\beta ad\left(  y\right)  \right)  \left(
z\right)  =\left[  \alpha x+\beta y,z\right]  =ad\left(  \alpha x+\beta
y\right)  \left(  z\right)  .
\]
Hence $ad\left(  \mathcal{G}\right)  $ is a subspace of $\mathcal{M\left(
\mathcal{G}\right)  }\subset\mathcal{L}\left(  \mathcal{G}\right)  .$

The map $ad$ is not surjective onto $\mathcal{L}\left(  \mathcal{G}\right)  $
i.e.
\begin{equation}
ad\left(  \mathcal{G}\right)  \subset\mathcal{M}\left(  \mathcal{G}\right)  .
\end{equation}


The compositions
\begin{equation}
\left(  ad\left(  x\right)  \circ ad\left(  y\right)  \right)  \left(
z\right)  =\left[  x,\left[  y,z\right]  \right]
\end{equation}
and
\begin{equation}
\left(  ad\left(  y\right)  \circ ad\left(  x\right)  \right)  \left(
z\right)  =\left[  y,\left[  x,z\right]  \right]
\end{equation}
produce%
\begin{equation}
\left(  ad\left(  x\right)  \circ ad\left(  y\right)  -ad\left(  y\right)
\circ ad\left(  x\right)  \right)  \left(  z\right)  =\left[  x,\left[
y,z\right]  \right]  -\left[  y,\left[  x,z\right]  \right]  ,
\end{equation}
while the Jacobi identity gives that
\begin{equation}
\left[  x,\left[  y,z\right]  \right]  -\left[  y,\left[  x,z\right]  \right]
=\left[  \left[  x,y\right]  ,z\right]  ,
\end{equation}
which shows that
\begin{equation}
\left[  \left[  x,y\right]  ,z\right]  =\left(  ad\left(  x\right)  \circ
ad\left(  y\right)  -ad\left(  y\right)  \circ ad\left(  x\right)  \right)
\left(  z\right)  .
\end{equation}
Hence the adjoint map of the product gives
\begin{equation}
ad\left(  \left[  x,y\right]  \right)  \left(  z\right)  =\left[  \left[
x,y\right]  ,z\right]
\end{equation}
and%
\begin{equation}
ad\left(  \left[  x,y\right]  \right)  =ad\left(  x\right)  \circ ad\left(
y\right)  -ad\left(  y\right)  \circ ad\left(  x\right)  =\left[  ad\left(
x\right)  ,ad\left(  y\right)  \right]  ,
\end{equation}
where we have used $\left[  ,\right]  $ to denote the lie product in both
$\mathcal{G}$ and $ad\left(  \mathcal{G}\right)  .$ Hence $ad$ is a Lie
algebra homomorphism and $ad\left(  \mathcal{G}\right)  $ is a Lie subalgebra
of $\mathcal{L}\left(  \mathcal{G}\right)  $ and a \emph{representation of}
$\mathcal{G}$ (the adjoint/regular representation)$.$

In particular this gives
\begin{equation}
ad\left(  \left[  0,x\right]  \right)  =ad\left(  0\right)  \circ ad\left(
x\right)  -ad\left(  x\right)  \circ ad\left(  0\right)  =0,\,\forall
x\in\mathcal{G}.
\end{equation}


If $\mathcal{G}$ is not simple, i.e contains non trivial ideals, which are
invariant subspaces to the action of the adjoint representation and the
adjoint representation decomposes as a direct sum of representations i.e
\begin{equation}
ad=%
%TCIMACRO{\tbigoplus \limits_{1\leq j\leq\kappa}}%
%BeginExpansion
{\textstyle\bigoplus\limits_{1\leq j\leq\kappa}}
%EndExpansion
ad_{j},
\end{equation}
where $\kappa$ is the number of ideals.

\begin{example}
Consider $\mathcal{G}=gl\left(  2,\mathbb{R}\right)  $ and the basis of
$gl\left(  2,\mathbb{R}\right)  $
\begin{equation}
\left\{  \left(
\begin{array}
[c]{cc}%
1 & 0\\
0 & 0
\end{array}
\right)  ,\left(
\begin{array}
[c]{cc}%
0 & 1\\
0 & 0
\end{array}
\right)  ,\left(
\begin{array}
[c]{cc}%
0 & 0\\
1 & 0
\end{array}
\right)  ,\left(
\begin{array}
[c]{cc}%
0 & 0\\
0 & 1
\end{array}
\right)  \right\}
\end{equation}
then the adjoint representation, contained in $\mathcal{L}\left(  gl\left(
2,\mathbb{R}\right)  \right)  $ wrt this basis is given by
\begin{align}
ad\left(
\begin{array}
[c]{cc}%
1 & 0\\
0 & 0
\end{array}
\right)   &  =\left(
\begin{array}
[c]{cccc}%
0 & 0 & 0 & 0\\
0 & 1 & 0 & 0\\
0 & 0 & -1 & 0\\
0 & 0 & 0 & 0
\end{array}
\right) \\
ad\left(
\begin{array}
[c]{cc}%
0 & 1\\
0 & 0
\end{array}
\right)   &  =\left(
\begin{array}
[c]{cccc}%
0 & 0 & 1 & 0\\
-1 & 0 & 0 & 0\\
0 & 0 & 0 & 0\\
0 & 0 & -1 & 0
\end{array}
\right) \\
ad\left(
\begin{array}
[c]{cc}%
0 & 0\\
1 & 0
\end{array}
\right)   &  =\left(
\begin{array}
[c]{cccc}%
0 & 0 & 0 & 0\\
0 & 0 & 0 & 0\\
1 & 0 & 0 & -1\\
0 & 1 & 0 & 0
\end{array}
\right) \\
ad\left(
\begin{array}
[c]{cc}%
0 & 0\\
0 & 1
\end{array}
\right)   &  =\left(
\begin{array}
[c]{cccc}%
0 & 0 & 0 & 0\\
0 & -1 & 0 & 0\\
0 & 0 & 1 & 0\\
0 & 0 & 0 & 0
\end{array}
\right)  ,
\end{align}
which is clearly a Lie subsalgebra of the Lie algebra $gl\left(
4,\mathbb{R}\right)  =M\left(  4,\mathbb{R}\right)  .$
\end{example}

\begin{example}
Again consider $\mathcal{G}=gl\left(  2,\mathbb{R}\right)  $ but with the
basis
\begin{equation}
\left\{  \frac{1}{2}\left(
\begin{array}
[c]{cc}%
1 & 0\\
0 & 1
\end{array}
\right)  ,\frac{1}{2}\left(
\begin{array}
[c]{cc}%
1 & 0\\
0 & -1
\end{array}
\right)  ,\left(
\begin{array}
[c]{cc}%
0 & 0\\
1 & 0
\end{array}
\right)  ,\left(
\begin{array}
[c]{cc}%
0 & 1\\
0 & 0
\end{array}
\right)  \right\}
\end{equation}
then
\begin{equation}
ad\left(
\begin{array}
[c]{cc}%
1 & 0\\
0 & 1
\end{array}
\right)  =\left(
\begin{array}
[c]{cccc}%
0 & 0 & 0 & 0\\
0 & 1 & 0 & 0\\
0 & 0 & -1 & 0\\
0 & 0 & 0 & 0
\end{array}
\right)  +\left(
\begin{array}
[c]{cccc}%
0 & 0 & 0 & 0\\
0 & -1 & 0 & 0\\
0 & 0 & 1 & 0\\
0 & 0 & 0 & 0
\end{array}
\right)  =\mathbb{O}_{4}%
\end{equation}
and
\begin{equation}
ad\left(
\begin{array}
[c]{cc}%
1 & 0\\
0 & 1
\end{array}
\right)  \left(  A\right)  =\mathbb{O}_{2}\,\forall A\in M\left(
2,\mathbb{R}\right)  .
\end{equation}
Thus in this case%
\begin{equation}
\operatorname{Ker}\left\{  ad\right\}  =\left\{  \left.  \left(
\begin{array}
[c]{cc}%
\lambda & 0\\
0 & \lambda
\end{array}
\right)  \right\vert \lambda\in\mathbb{R}\right\}  =gl\left(  1,\mathbb{R}%
\right)
\end{equation}
and
\begin{equation}
gl\left(  2,\mathbb{R}\right)  =gl\left(  1,\mathbb{R}\right)  \oplus
sl\left(  2,\mathbb{R}\right)  .
\end{equation}
Observe that $gl\left(  2,\mathbb{R}\right)  $ is not semisimple as $gl\left(
1,\mathbb{R}\right)  $ is abelian (see later).
\end{example}

\begin{theorem}
Let $G$ be a Lie group. Denote by $G_{0}$ the connected component of unity.
Then $G_{0}$ is a \emph{normal subgroup} of $G$ and is a Lie group itself. The
quotient group $G/G_{0}$ is discrete.
\end{theorem}

\begin{remark}
This theorem mostly reduces the study of \emph{arbitrary Lie groups} to the
s\emph{tudy of finite groups and connected Lie groups}. In fact, one can go
further and \emph{reduce the study of \underline{\emph{connected Lie}} groups
to connected \underline{\emph{simply-connected }}Lie groups}.
\end{remark}

\begin{remark}
$SO\left(  r,\mathbb{R}\right)  $ \emph{is connected but \underline{\emph{NOT}%
} simply connected. On the other hand }$SU\left(  r\right)  $ is \emph{simply
connected. Both these groups have the same Lie algebra, which is the same Lie
algebra of the group }$O\left(  r,\mathbb{R}\right)  $ i.e.
\begin{equation}
so\left(  r,\mathbb{R}\right)  =o\left(  r,\mathbb{R}\right)  =su\left(
r\right)  .
\end{equation}

\end{remark}

\begin{theorem}
If $G$ is a connected Lie group then its universal cover $\mathcal{\rho
}\left(  G\right)  $ has a canonical structure of a Lie group such that the
covering map%
\begin{equation}
\mathcal{\rho}:G\rightarrow\mathcal{\rho}\left(  G\right)
\end{equation}
is a morphism of Lie groups, and
\begin{equation}
\ker\left\{  \mathcal{\rho}\right\}  =\Pi_{1}\left(  G\right)  \subset
\mathcal{Z}\left(  G\right)
\end{equation}
is a group. Morever, in this case $\Pi_{1}\left(  G\right)  $ is a discrete
central subgroup in $\mathcal{\rho}\left(  G\right)  $.
\end{theorem}

\paragraph{Lie Group Adjoint Representation}

As $\mathcal{G}$ is a vector space one can consider the general linear matrix
group $GL\left(  \mathcal{G}\right)  =\left\{  \left.  g\right\vert
g\in\mathcal{L}\left(  \mathcal{G}\right)  \ni\exists g^{-1}\right\}  .$ One
can define a representation, $Ad,$
\begin{equation}
Ad:GL\left(  \mathcal{G}\right)  \rightarrow Ad\left(  GL\left(
\mathcal{G}\right)  \right)  \subset GL\left(  \mathcal{L}\left(
\mathcal{G}\right)  \right)  \subset\mathcal{L}\left(  \mathcal{L}\left(
\mathcal{G}\right)  \right)
\end{equation}
of $GL\left(  \mathcal{G}\right)  $ by
\begin{align}
Ad\left(  g\right)   &  :\mathcal{L}\left(  \mathcal{G}\right)  \rightarrow
\mathcal{L}\left(  \mathcal{G}\right) \\
\left(  Ad\left(  g\right)  \right)  \left(  X\right)   &  =gXg^{-1},
\end{align}%
\begin{align}
Ad\left(  g\right)  \circ Ad\left(  h\right)   &  =Ad\left(  gh\right) \\
Ad\left(  g^{-1}\right)   &  =\left(  Ad\left(  g\right)  \right)  ^{-1}\\
Ad\left(  g\right)  \circ\left(  Ad\left(  g\right)  \right)  ^{-1}  &
=I_{\mathcal{G}}.
\end{align}
The elements of $Ad\left(  GL\left(  \mathcal{G}\right)  \right)  $ are Lie
algebra mapping as
\begin{equation}
\left(  Ad\left(  g\right)  \right)  \left(  \left[  X,Y\right]  \right)
=\left[  \left(  Ad\left(  g\right)  \right)  \left(  X\right)  ,\left(
Ad\left(  g\right)  \right)  \left(  Y\right)  \right]  .
\end{equation}
The restriction of $Ad\left(  GL\left(  \mathcal{G}\right)  \right)  $ to
$GL\left(  \mathcal{G}\right)  $ gives group conjugations
\begin{equation}
Ad\left(  g\right)  \left(  h\right)  =ghg^{-1},
\end{equation}
which belong to $\operatorname*{Inn}\left\{  GL\left(  \mathcal{G}\right)
\right\}  ,$ which is a normal subgroup of $\operatorname{Aut}\left\{
GL\left(  \mathcal{G}\right)  \right\}  $ showing that
\begin{equation}
\operatorname*{Out}\left\{  GL\left(  \mathcal{G}\right)  \right\}
=\frac{\operatorname{Aut}\left\{  GL\left(  \mathcal{G}\right)  \right\}
}{\operatorname*{Inn}\left\{  GL\left(  \mathcal{G}\right)  \right\}  }%
\end{equation}
is a group. One can show that
\begin{equation}
\operatorname*{Inn}\left\{  GL\left(  \mathcal{G}\right)  \right\}  \cong%
\frac{GL\left(  \mathcal{G}\right)  }{\mathcal{Z}\left(  GL\left(
\mathcal{G}\right)  \right)  },
\end{equation}
thus if
\begin{equation}
\mathcal{Z}\left(  GL\left(  \mathcal{G}\right)  \right)  =I_{\mathcal{G}}%
\end{equation}
then
\begin{equation}
\operatorname*{Inn}\left\{  GL\left(  \mathcal{G}\right)  \right\}  \cong
GL\left(  \mathcal{G}\right)  .
\end{equation}


\begin{remark}
Let $G$ be a \emph{compact simple simply connected Lie group}. Then any
automorphism of $G$ determines an automorphism of its Lie algebra
$\mathcal{G}$ and visa versa. So $Aut(G)$ is naturally isomorphic to the
linear group $Aut(\mathcal{G})$ and one has
\begin{equation}
\operatorname{Aut}\left\{  G\right\}  \cong\operatorname{Aut}\left\{
\mathcal{G}\right\}  \cong\operatorname*{Inn}\left\{  \mathcal{G}\right\}
\rtimes\operatorname*{Out}\left\{  G\right\}  \cong\operatorname*{Inn}\left\{
\mathcal{G}\right\}  \rtimes\operatorname{Aut}\left\{  \mathcal{D}%
_{\mathcal{G}}\right\}  .
\end{equation}

\end{remark}

\begin{remark}
if $G$ is any connected Lie group, the \emph{adjoint representation} of $G$ on
its Lie algebra $\mathcal{G}$ has kernel precisely the center of $G$, and so
furnishes a faithful linear representation of $G/\mathcal{Z}(G)$. \emph{All of
the projective linear groups} arise in this way.
\end{remark}

\begin{example}%
\begin{equation}
\operatorname{Aut}\left\{  SU\left(  n\right)  \right\}  \cong\left\{
\begin{array}
[c]{c}%
PSU\left(  n\right)  \rtimes\frac{\mathbb{Z}}{2\mathbb{Z}};n\geq3\\
PU\left(  2\right)  ;n=2.
\end{array}
\right.
\end{equation}

\end{example}

\begin{remark}
If $\mathbb{F}$ is any field, the quotient group $GL(n,\mathbb{F}%
)/\mathcal{Z}\left(  GL(n,\mathbb{F})\right)  $ is called the \emph{projective
group} and is written $PGL(n,\mathbb{F})$. The elements of this group are
called projective transformations. The centre $\mathcal{Z}\left(
GL(n,\mathbb{F})\right)  =$ scalar transformations $\left\{  \left.  \lambda
I_{n}\right\vert \lambda\in\mathbb{F}\right\}  .$
\end{remark}

\begin{remark}
The exponential is defined in the associative algebra $\mathcal{L}\left(
\mathcal{G}\right)  $
\begin{equation}
\exp:\mathcal{L}\left(  \mathcal{G}\right)  \rightarrow\mathcal{L}\left(
\mathcal{G}\right)
\end{equation}
by the power series
\begin{equation}
\exp\left\{  A\right\}  =\sum_{1\leq n\leq\infty}\frac{1}{n!}A^{n}.
\end{equation}
If $\mathcal{L}\left(  \mathcal{G}\right)  $ is a Banach algebra then $\exp$
converges everywhere and is an entire function (analytic everywhere).
\end{remark}

\paragraph{Adjoint Representation, $Ad\left(  GL\left(  \mathcal{G}\right)
\right)  ,$ of in terms of the adjoint representation, $ad\left(
\mathcal{G}\right)  ,$ of $\mathcal{G}.$}

The exponential maps $\mathcal{L}\left(  \mathcal{G}\right)  $ into $GL\left(
\mathcal{G}\right)  $ i.e.
\begin{align}
\exp &  :\mathcal{L}\left(  \mathcal{G}\right)  \rightarrow GL\left(
\mathcal{G}\right) \\
g_{A}  &  =\exp\left\{  A\right\}  ;\,A\in\mathcal{L}\left(  \mathcal{G}%
\right)  ,
\end{align}
as $\exp\left\{  A\right\}  \forall A\in\mathcal{L}\left(  \mathcal{G}\right)
$ is always invertible i.e.
\begin{equation}
\exp\left\{  A\right\}  \exp\left\{  -A\right\}  =I_{\mathcal{G}}.
\end{equation}


\begin{remark}
$\exp$ is \emph{not necessarily injective} e.g. consider Lie algebra
$so\left(  2,\mathbb{R}\right)  $ given by
\begin{equation}
so\left(  2,\mathbb{R}\right)  =\left\{  \left.  A\in M\left(  2,\mathbb{R}%
\right)  \right\vert A=-A^{t}\right\}
\end{equation}
and the action of the expontial map on $so\left(  2,\mathbb{R}\right)  $ i.e.
\begin{equation}
\exp\left\{  \left(
\begin{array}
[c]{cc}%
0 & \theta\\
-\theta & 0
\end{array}
\right)  \right\}  =\left(
\begin{array}
[c]{cc}%
\cos\theta & \sin\theta\\
-\sin\theta & \cos\theta
\end{array}
\right)  ,
\end{equation}
which shows that $\theta+2n\pi$ gives the same matrix under the exponential
map. In particular this shows that
\begin{equation}
\log\left\{  \left(
\begin{array}
[c]{cc}%
\cos\theta & \sin\theta\\
-\sin\theta & \cos\theta
\end{array}
\right)  \right\}  =\left(
\begin{array}
[c]{cc}%
0 & \theta+2n\pi\\
-\left(  \theta+2n\pi\right)  & 0
\end{array}
\right)
\end{equation}
and $\left(
\begin{array}
[c]{cc}%
\cos\theta & \sin\theta\\
-\sin\theta & \cos\theta
\end{array}
\right)  $ has infinitely many logarithms.

One can observe that the action of the exponential map on $so\left(
2,\mathbb{R}\right)  $ produces rotations of $\mathbb{R}^{2}$ that have
determinant $+1$, and form the special orthogonal group
\begin{equation}
SO\left(  2,\mathbb{R}\right)  =\left\{  \left.  T\in M\left(  2,\mathbb{R}%
\right)  \right\vert TT^{t}=T^{t}T=I_{2};\det\left\{  A\right\}  =1\right\}  .
\end{equation}

\end{remark}

\begin{remark}
$\exp$ is \emph{not necessarily surjective} e.g. consider reflections $T\notin
SO\left(  2,\mathbb{R}\right)  ,$ where
\begin{equation}
TT^{t}=T^{t}T=I_{2}%
\end{equation}
and%
\begin{equation}
\det\left\{  T\right\}  =\det\left\{  \left(
\begin{array}
[c]{cc}%
-\cos\theta & \sin\theta\\
\sin\theta & \cos\theta
\end{array}
\right)  \right\}  =-1.
\end{equation}
As
\begin{equation}
\det\left\{  \exp\left\{  A\right\}  \right\}  =\exp\left\{  \operatorname{Tr}%
\left\{  A\right\}  \right\}  >0
\end{equation}
$T$ can never be in the image of $\exp,$ thus
\begin{equation}
\exp:so\left(  2,\mathbb{R}\right)  \rightarrow O\left(  2,\mathbb{R}\right)
\end{equation}
is not surjective.
\end{remark}

\begin{remark}
The orthogonal group $O\left(  2,\mathbb{R}\right)  $ is given by the union of
two disconnected components (rotations and reflections)
\begin{equation}
O\left(  2,\mathbb{R}\right)  =\left\{  \left.  \left(
\begin{array}
[c]{cc}%
-\cos\theta & \sin\theta\\
\sin\theta & \cos\theta
\end{array}
\right)  ,\left(
\begin{array}
[c]{cc}%
\cos\theta & \sin\theta\\
-\sin\theta & \cos\theta
\end{array}
\right)  \right\vert \theta\in\left[  0,2\pi\right)  \right\}  .
\end{equation}

\end{remark}

\begin{remark}
The name rotation matrices for the elements of $SO\left(  r,\mathbb{R}\right)
$ is justified by

\begin{proposition}
.Any matrix in $SO\left(  r,\mathbb{R}\right)  $ is \emph{orthogonally
similar} to a block diagonal of the form%
\begin{equation}
A_{1}\oplus A_{2}\oplus\text{\textperiodcentered\textperiodcentered
\textperiodcentered}\ \oplus A_{r}%
\end{equation}
,where each $A_{j}$ is $\left\{  1\right\}  $ or a $2\times2$ rotation matrix
of the type%
\begin{equation}
\left(
\begin{array}
[c]{cc}%
\cos\theta & \sin\theta\\
-\sin\theta & \cos\theta
\end{array}
\right)  .
\end{equation}
for some $\U{3b8} \in\mathbb{R}$.
\end{proposition}
\end{remark}

\begin{lemma}
It can be shown that if $A,B\in\mathcal{L}\left(  \mathcal{G}\right)  $ and
$A$ is invertible then we can write
\begin{equation}
Ad\left(  g_{A}\right)  \left(  B\right)  =g_{A}Bg_{A}^{-1}=\exp\left\{
A\right\}  B\exp\left\{  -A\right\}
\end{equation}
and one can show that
\begin{equation}
Ad\left(  g_{A}\right)  \left(  B\right)  =\exp\left\{  ad\left(  A\right)
\right\}  \left(  B\right)  .
\end{equation}
Thus $Ad\left(  g_{A}\right)  $ is an \emph{inner automorphism} of
$\mathcal{L}\left(  \mathcal{G}\right)  $ i.e. $Ad\left(  g_{A}\right)  \in$
$\operatorname*{Inn}\left\{  \mathcal{L}\left(  \mathcal{G}\right)  \right\}
\subset GL\left(  \mathcal{L}\left(  \mathcal{G}\right)  \right)
\subset\mathcal{L}\left(  \mathcal{L}\left(  \mathcal{G}\right)  \right)  $
and $ad\left(  A\right)  \in\operatorname{Der}\left\{  \mathcal{L}\left(
\mathcal{G}\right)  \right\}  ,$ the $\emph{Derivations}$ of $\mathcal{L}%
\left(  \mathcal{G}\right)  \subset\mathcal{L}\left(  \mathcal{L}\left(
\mathcal{G}\right)  \right)  .$
\end{lemma}

\emph{Automorphisms} of $\mathcal{L}\left(  \mathcal{G}\right)  $ can be inner
or outer, where
\begin{equation}
\operatorname*{Out}\left\{  \mathcal{L}\left(  \mathcal{G}\right)  \right\}
=\frac{\operatorname{Aut}\left\{  \mathcal{L}\left(  \mathcal{G}\right)
\right\}  }{\operatorname*{Inn}\left\{  \mathcal{L}\left(  \mathcal{G}\right)
\right\}  }.
\end{equation}
\emph{Derivations} can be inner or outer as well, all inner derivations are of
the form $ad\left(  A\right)  ,$ where $A\in\mathcal{L}\left(  \mathcal{G}%
\right)  $.

If $G$ is a real or complex Lie group with a \emph{semi-simple} Lie algebra,
then the inner automorphisms constitute precisely the identity component of
the group $GL\left(  \mathcal{G}\right)  $ of automorphisms of $\mathcal{G}$.

One can observe that
\begin{equation}
\mathcal{L}\left(  \mathcal{G}\right)  \equiv gl\left(  GL\left(
\mathcal{G}\right)  \right)
\end{equation}
and the adjoint representation of $GL\left(  \mathcal{G}\right)  $ is given
by
\begin{equation}
Ad\left(  GL\left(  \mathcal{G}\right)  \right)  =\left\{  \left.
\exp\left\{  ad\left(  x\right)  \right\}  \right\vert x\in\mathcal{L}\left(
\mathcal{G}\right)  \equiv gl\left(  GL\left(  \mathcal{G}\right)  \right)
\right\}  .
\end{equation}
The identity element of $Ad\left(  G\right)  $ is $I_{\mathcal{G}}$ and
\begin{align}
Ad\left(  g\right)   &  =\exp\left\{  ad\left(  x\right)  \right\} \\
\left(  Ad\left(  g\right)  \right)  ^{-1}  &  =\exp\left\{  -ad\left(
x\right)  \right\} \\
Ad\left(  g\right)  \left(  Ad\left(  g\right)  \right)  ^{-1}  &
=\exp\left\{  ad\left(  x\right)  \right\}  \exp\left\{  -ad\left(  x\right)
\right\}  =\exp\left\{  ad\left(  x\right)  -ad\left(  x\right)  \right\}
=\exp\left\{  ad\left(  0\right)  \right\}  =\exp\left\{  0\right\}
=I_{\mathcal{G}}.
\end{align}
The group product in the $Ad$ representation is given by
\begin{equation}
\exp\left\{  ad\left(  x\right)  \right\}  \exp\left\{  ad\left(  y\right)
\right\}  =\exp\left\{  \mathrm{CBH}\left(  ad\left(  x\right)  ,ad\left(
x\right)  \right)  \right\}  ,
\end{equation}
where $\mathrm{CBH}$ is the \emph{Cambell Baker Hausdorff }function, whose
arguments are operators or elements of a Lie algebra. It is defined by a
formal power series in the Lie algebra $\mathcal{G}$ by \emph{ }%
\begin{align}
\mathrm{CBH}  &  \mathrm{:}\mathcal{G}\times\mathcal{G}\rightarrow
\mathcal{G}\\
\mathrm{CBH}\left(  x,y\right)   &  =\log\left\{  e^{x}e^{y}\right\}
=x+y+\frac{1}{2}\left[  x,y\right]  +\frac{1}{12}\left[  x,\left[  x,y\right]
\right]  +\frac{1}{12}\left[  y,\left[  x,y\right]  \right]  +\cdots,
\end{align}
for the linear operators, $ad\left(  \mathcal{G}\right)  ,$ acting in
$\mathcal{G}$ as
\begin{align}
\mathrm{CBH}  &  \mathrm{:}ad\left(  \mathcal{G}\right)  \times ad\left(
\mathcal{G}\right)  \rightarrow ad\left(  \mathcal{G}\right) \\
\mathrm{CBH}\left(  x,y\right)   &  =\log\left\{  e^{ad\left(  x\right)
}e^{ad\left(  x\right)  }\right\} \\
&  =ad\left(  x\right)  +ad\left(  y\right)  +\frac{1}{2}\left[  ad\left(
x\right)  ,ad\left(  y\right)  \right]  +\frac{1}{12}\left[  ad\left(
x\right)  ,\left[  ad\left(  x\right)  ,ad\left(  y\right)  \right]  \right]
+\frac{1}{12}\left[  ad\left(  y\right)  ,\left[  ad\left(  x\right)
,ad\left(  y\right)  \right]  \right]  +\cdots
\end{align}
and for all linear operators in $\mathcal{L}\left(  \mathcal{G}\right)  ,$
acting in $\mathcal{G}$ as
\begin{align}
\mathrm{CBH}  &  \mathrm{:}\mathcal{L}\left(  \mathcal{G}\right)
\times\mathcal{L}\left(  \mathcal{G}\right)  \rightarrow\mathcal{L}\left(
\mathcal{G}\right) \\
\mathrm{CBH}\left(  A,B\right)   &  =\log\left\{  e^{A}e^{B}\right\} \\
&  =A+B+\frac{1}{2}\left[  A,B\right]  +\frac{1}{12}\left[  A,\left[
A,B\right]  \right]  +\frac{1}{12}\left[  B,\left[  A,B\right]  \right]
+\cdots.
\end{align}
It can be shown that if $A,B\in\mathcal{L}\left(  \mathcal{G}\right)  $ and
$A$ is invertible then we can write
\begin{equation}
\left(  Ad\left(  A\right)  \right)  \left(  B\right)  =ABA^{-1}%
\end{equation}
and one can show that
\begin{equation}
\left(  Ad\left(  e^{X}\right)  \right)  \left(  B\right)  =e^{X}Be^{-X}%
=\exp\left\{  ad\left(  X\right)  \right\}  \left(  B\right)  ,
\end{equation}
where
\begin{equation}
A=e^{X}=e^{\log\left\{  A\right\}  }.
\end{equation}
It is not always true that
\begin{equation}
\mathrm{CBH:}ad\left(  \mathcal{G}\right)  \times ad\left(  \mathcal{G}%
\right)  \rightarrow ad\left(  \mathcal{G}\right)  .
\end{equation}


The group $Ad\left(  G\right)  $ acts in $\mathcal{G}$ by
\begin{align}
\exp\left\{  ad\left(  x\right)  \right\}  \left(  y\right)   &  =y+ad\left(
x\right)  \left(  y\right)  +\frac{1}{2}\left(  ad\left(  x\right)  \circ
ad\left(  x\right)  \right)  \left(  y\right)  +\frac{1}{3!}\left(  ad\left(
x\right)  \circ ad\left(  x\right)  \circ ad\left(  x\right)  \right)  \left(
y\right)  +\cdots\\
&  =y+\left[  x,y\right]  +\frac{1}{2}\left[  x,\left[  x,y\right]  \right]
+\frac{1}{3!}\left[  x,\left[  x,\left[  x,y\right]  \right]  \right]  +\cdots
\end{align}


\subsection{SemiSimple Lie Algebras}

A Lie algebra is \emph{semisimple} if it is a direct sum of \emph{simple Lie
algebras, i.e., \underline{\emph{non-abelian}} Lie algebras }$\mathcal{G}%
$\emph{ whose only ideals are }$\{0\}$\emph{ and }$\mathcal{G}$\emph{ itself.}
It is important to \emph{emphasize that a one-dimensional Lie algebra} (which
is necessarily abelian) is by definition \emph{not considered a simple Lie
algebra}, even though such an algebra certainly has no nontrivial ideals.
\emph{Thus, one-dimensional algebras are not allowed as summands in a
semisimple Lie algebra}.

Unless otherwise stated, $\mathcal{G}$ is a non-zero finite-dimensional Lie
algebra over a field of characteristic $0$. The following conditions are equivalent:

\begin{enumerate}
\item $\mathcal{G}$ is semisimple

\item the \emph{Killing form} on the Lie algebra is, $K\left(  x,y\right)
=\operatorname{Tr}\left\{  ad\left(  x\right)  ad\left(  y\right)  \right\}  $
is non-degenerate,

\item $\mathcal{G}$ has no non-zero abelian ideals,

\item $\mathcal{G}$ has no non-zero \emph{solvable} ideals,

$\lceil\mathcal{G}$ is solvable if
\begin{equation}
\left[  \mathcal{G},\mathcal{G}\right]  \supseteq\left[  \left[
\mathcal{G},\mathcal{G}\right]  ,\left[  \mathcal{G},\mathcal{G}\right]
\right]  \supseteq\left[  \left[  \left[  \mathcal{G},\mathcal{G}\right]
,\left[  \mathcal{G},\mathcal{G}\right]  \right]  ,\left[  \left[
\mathcal{G},\mathcal{G}\right]  ,\left[  \mathcal{G},\mathcal{G}\right]
\right]  \right]  \supseteq\cdots\rightarrow0
\end{equation}
$\rfloor$

\item The \emph{radical} (maximal solvable ideal) of $\mathcal{G}$ is zero.
\end{enumerate}

\paragraph{Examples}

Semisimple Lie algebras over $\mathbb{C}$ are classified by the root system
associated to their \emph{Cartan} subalgebras

$\lceil$A Cartan subalgebra is a subalgebra that is nilpotent
\begin{equation}
\exists n\ni\left(
%TCIMACRO{\tprod \limits_{1\leq j\leq n}}%
%BeginExpansion
{\textstyle\prod\limits_{1\leq j\leq n}}
%EndExpansion
ad\left(  X_{j}\right)  \right)  \left(  Y\right)  =0\,\forall\left\{
X_{1},\cdots,X_{n},Y\right\}  \in\mathcal{L}%
\end{equation}
and self normalizing
\begin{equation}
\left[  X,Y\right]  \in\mathcal{L\,\,\,}\forall X\in\mathcal{L}\Rightarrow
Y\in\mathcal{L}.
\end{equation}
$\rfloor$

and their \emph{root systems} are classified by their \emph{Dynkin diagrams}.
Examples of semisimple Lie algebras, with notation coming from their Dynkin
diagrams, are:%

\begin{align}
A_{n}  &  :sl\left(  n+1,\mathbb{C}\right)  ;\text{ The special linear Lie
algebra, }n>0\\
B_{n}  &  :so\left(  2n+1,\mathbb{C}\right)  ;\text{ The special orthogonal
Lie algebra of odd dimension, }n>0\\
C_{n}  &  :sp\left(  n+1,\mathbb{C}\right)  ;\text{ The symplectic Lie
algebra, }n>0\\
D_{n}  &  :so\left(  2n,\mathbb{C}\right)  ,n>1;\text{ The special orthogonal
Lie algebra of even dimension}%
\end{align}
($so\left(  2,\mathbb{C}\right)  $ is one-dimensional and commutative and
therefore not semisimple).

These Lie algebras are numbered so that $n$ is their rank. Almost all of these
semisimple Lie algebras are actually simple. These four families, together
with five exceptions ($E_{6},E_{7},E_{8},F_{4},$ and $G_{2}$), are in fact the
only simple Lie algebras over the complex numbers.

\paragraph{Classification}

The simple Lie algebras are classified by connected Dynkin diagrams%

%TCIMACRO{\FRAME{ftbpF}{3.6106in}{1.2835in}{0pt}{}{}{index.png}%
%{\special{ language "Scientific Word";  type "GRAPHIC";
%maintain-aspect-ratio TRUE;  display "USEDEF";  valid_file "F";
%width 3.6106in;  height 1.2835in;  depth 0pt;  original-width 1.7916in;
%original-height 0.6276in;  cropleft "0";  croptop "1";  cropright "1";
%cropbottom "0"; }}}%
%BeginExpansion
\begin{figure}[ptb]%
\centering
\includegraphics[
natheight=0.627600in,
natwidth=1.791600in,
height=1.2835in,
width=3.6106in
]%
{}%
\end{figure}
%EndExpansion


Every semisimple Lie algebra over an algebraically closed field of
characteristic $0$ is a direct sum of simple Lie algebras (by definition), and
the finite-dimensional simple Lie algebras fall in four families
\begin{equation}
A_{n},B_{n},C_{n},D_{n}%
\end{equation}
with five exceptions
\begin{equation}
E_{6},E_{7},E_{8},F_{4},\text{and }G_{2}.
\end{equation}
Simple Lie algebras are classified by the connected Dynkin diagrams, while
semisimple Lie algebras correspond to Dynkin diagrams (not necessarily
connected), where each component of the diagram corresponds to a summand of
the decomposition of the semisimple Lie algebra into simple Lie algebras.

The classification proceeds by considering a \emph{Cartan} subalgebra and the
adjoint action of the Lie algebra on this subalgebra. The \emph{root system}
of the action determines the original Lie algebra and has a very constrained
form, which can be classified by the Dynkin diagrams.

The classification is widely considered one of the most elegant results in
mathematics -- a brief list of axioms yields, via a relatively short proof, a
complete but non-trivial classification with surprising structure. This should
be compared to the classification of finite simple groups, which is
significantly more complicated.

The enumeration of the four families is non-redundant and consists only of
\emph{simple algebras} if $n\geq1$ for $A_{n};n\geq2$ for $B_{n};n\geq3$ for
$C_{n}$ and $n\geq4$ for $D_{n}$. If one starts numbering lower, the
enumeration is redundant, and one has exceptional isomorphisms between simple
Lie algebras, which are reflected in isomorphisms of Dynkin diagrams; the
$E_{n}$ can also be extended down, but below $E_{6}$ they are isomorphic to
other, non-exceptional algebras.

\emph{Over a non-algebraically closed field, e.g $\mathbb{R}$, the
classification is more complicated} -- one classifies simple Lie algebras over
the algebraic closure, then for each of these, one classifies simple Lie
algebras over the original field which have this form (over the closure). For
example, to classify simple real Lie algebras, one classifies real Lie
algebras with a given complexification, which are known as real forms of the
complex Lie algebra; this can be done by \emph{Satake diagrams}, which are
Dynkin diagrams with additional data ("decorations").

\subsubsection{Compact Lie groups}

\begin{definition}
A compact (topological) group is a topological group whose topology is compact.
\end{definition}

\begin{definition}
A topological space is \emph{compact} $\operatorname*{Iff}$ every open cover
has a finite open subcover.
\end{definition}

\begin{remark}
Compactness generalizes the concepts of closed and bounded.
\end{remark}

\begin{remark}
Compact groups are a natural generalization of finite groups with the discrete
topology and have properties that carry over in significant fashion.
\end{remark}

There are two definitions of a compact Lie algebra. Extrinsically and topologically:

\begin{definition}
A compact Lie algebra is the Lie algebra of a compact Lie group;
\end{definition}

\begin{remark}
This definition includes tori.
\end{remark}

\begin{definition}
Intrinsically and algebraically, a compact Lie algebra is a real Lie algebra
whose Killing form is negative definite
\end{definition}

\begin{remark}
This definition is more restrictive and excludes tori. .
\end{remark}

\begin{lemma}
Let $K$ be a \emph{compact real }Lie group with Lie algebra $\mathcal{K}$.
Then the complexification
\begin{equation}
\mathcal{G}=\mathcal{K}\oplus i\mathcal{K}=\mathbb{C\otimes}\mathcal{K}%
\end{equation}
of $\mathcal{K}$ is\emph{ reductive}, that is, \emph{the direct sum }of a
complex semisimple Lie algebra and a commutative algebra.
\end{lemma}

\begin{lemma}
If $K$ is \emph{simply connected}, then $\mathcal{G}$ is \emph{semisimple}.

\begin{corollary}
Conversely, every complex semisimple Lie algebra $\mathcal{G}$ has a compact
real form $\mathcal{K}$ where $\mathcal{K}$ is the Lie algebra of a simply
connected compact Lie group.
\end{corollary}
\end{lemma}

\begin{remark}
For example, the complex semisimple Lie algebra $sl\left(  n,\mathbb{C}%
\right)  $ is the complexification of $su\left(  n\right)  $ the Lie algebra
of the simply connected compact group $SU(n)$.
\end{remark}

\begin{remark}
It is possible to develop the theory of complex semisimple Lie algebras from
the compact group perspective, leading to a simpler way to develop the
existence and properties of Cartan subalgebras.
\end{remark}

\begin{lemma}
A compact Lie algebra is the smallest real form of a corresponding complex Lie
algebra, namely the complexification.
\end{lemma}

\paragraph{History}

The semisimple Lie algebras over the complex numbers were first classified by
\emph{Wilhelm Killing (1888--90)}, though his proof lacked rigor. His proof
was made rigorous by \emph{\'{E}lie Cartan (1894)} in his Ph.D. thesis, who
also classified semisimple real Lie algebras. This was subsequently refined,
and the present classification by Dynkin diagrams was given by then
\emph{22-year-old Eugene Dynkin in 1947}. Some minor modifications have been
made (notably by J. P. Serre), but the proof is unchanged in its essentials
and can be found in any standard reference, such as (Humphreys 1972).

\paragraph{Properties}

\subparagraph{Complete reducibility}

A consequence of semisimplicity is a theorem due to Weyl: \emph{every
\underline{\emph{finite-dimensional}} representation is completely reducible};
that is for every invariant subspace of the representation there is an
invariant complement. \emph{Infinite-dimensional representations of semisimple
Lie algebras are not in general completely reducible.}

\subparagraph{Centerless}

Since the center of a Lie algebra $\mathcal{G}$ is an \emph{abelian ideal},
then if $\mathcal{G}$ is semisimple then its centre is zero. ($gl\left(
n,\mathbb{C}\right)  $ has a non-trivial center, as it is \emph{not
}semisimple.) In other words, the \emph{adjoint representation }$ad$ of a
semisimple Lie algebra
\begin{equation}
ad:\mathcal{G}\rightarrow\mathcal{L}\left(  \mathcal{G}\right)
\end{equation}
is \emph{injective i.e faithful}.

Moreover, it can be shown that the dimension of the Lie algebra,
$\operatorname{Der}\left(  \mathcal{G}\right)  ,$ of derivations on
$\mathcal{G}$ is equal to the dimension of $\mathcal{G}$. Hence, $\mathcal{G}$
is \emph{Lie algebra isomorphic} to $\operatorname{Der}\left(  \mathcal{G}%
\right)  .$ (This is a special case of Whitehead's lemma.) Every ideal,
quotient and product of semisimple Lie algebras is again semisimple.

\subparagraph{Linear}

The adjoint representation of semisimple algebras is injective, and so a
semisimple Lie algebra is also a linear Lie algebra under the adjoint
representation. This may lead to some ambiguity, as every Lie algebra is
already linear with respect to some other vector space (Ado's theorem),
although not necessarily via the adjoint representation. But in practice, such
ambiguity rarely occurs.

\subparagraph{Jordan decomposition}

Any \emph{endomorphism}, $X,$ of a finite-dimensional vector space over an
algebraically closed field can be decomposed uniquely into a
\emph{diagonalizable (or semisimple) and nilpotent part}%

\begin{equation}
X=S+N
\end{equation}
such that $S$ and $N$ commute with each other. Moreover, each of $S$ and $N$
is a polynomial in $X$, which is a consequence of the \emph{Jordan
decomposition.}

If $X\in\mathcal{G}$ then the image of $X$ under the adjoint map decomposes
as
\begin{equation}
ad\left(  X\right)  =ad\left(  S\right)  +ad\left(  N\right)  .
\end{equation}


The elements $S$ and $N$ are unique elements of $\mathcal{G}$ such that $N$ is
nilpotent, $S$ is semisimple, $S$ and $N$ commute, and for which such a
decomposition holds. This \emph{abstract Jordan decomposition} factors through
any representation $\Pi$ i.e.
\begin{equation}
\Pi\left(  X\right)  =\Pi\left(  S\right)  +\Pi\left(  N\right)  .
\end{equation}


\subparagraph{Rank}

The rank of a complex semisimple Lie algebra is the dimension of any of its
\emph{Cartan subalgebra}.

\subsection{Representation Theory}

\begin{theorem}
\emph{Weyl's theorem:} Complete reducibility implies that every
finite-dimensional representation of a semisimple Lie algebra decomposes as a
direct sum of irreducible representations. The finite-dimensional, irreducible
representations, meanwhile, are classified by a theorem of the highest weight.
\end{theorem}

\subsubsection{Cartan subalgebras and root systems}

\paragraph{Definition of a Cartan subalgebra}

Although there is a theory of Cartan subalgebras for any Lie algebra, the
concept has particular importance and in a special form for the case of
\emph{complex semisimple Lie algebras.}

\begin{definition}
If $\mathcal{G}$ is a complex semisimple Lie algebra, we say that
$\mathcal{H}$ is a\emph{ Cartan subalgebra} if $\mathcal{H}$ is a
\emph{maximal commutative} subalgebra of $\mathcal{G}$ and if $ad\left(
h\right)  $ is diagonalizable for each $h\in\mathcal{H}$.
\end{definition}

\begin{remark}
An important first step in the study of semisimple Lie algebras is to prove
the existence of Cartan subalgebras, and their uniqueness up to automorphism.
(If one assumes the existence of a compact real form, the existence of a
Cartan subalgebra is much simpler to establish. In that case, $\mathcal{H}$
may be taken as the complexification of the Lie algebra of a\emph{ maximal
torus} of the compact group.)
\end{remark}

\begin{definition}
$\bigskip$A torus \emph{in a compact Lie group} $\mathcal{G}$ is a compact,
connected, abelian Lie subgroup of $\mathcal{G}$ and therefore isomorphic to
the standard torus
\begin{equation}
T_{n}=\overset{n}{\overbrace{\,S^{1}\times S^{1}\times\cdots\times S^{1}}}).
\end{equation}
A maximal torus is one which is maximal among such subgroups
\end{definition}

\begin{remark}
Since Cartan subalgebras of a semisimple Lie algebra $\mathcal{G}$ are unique
up to automorphisms of $\mathcal{G}$, all Cartan subalgebras have the same
dimension. This common dimension is the rank of $\mathcal{G}$.
\end{remark}

\paragraph{Root systems}

Given a Cartan subalgebra $\mathcal{H}$ of $\mathcal{G}$, one defines \emph{a
root} to be a nonzero element $\U{3b1} $ of $\mathcal{H}^{d}$ for which there
exists a nonzero $X\in\mathcal{G}$ such that
\begin{equation}
ad\left(  h\right)  \left(  X\right)  =\left[  h,X\right]  =\alpha\left(
X\right)  \quad\forall h\in\mathcal{H}.
\end{equation}
That is to say, the roots are the \emph{nonzero weights} $\lceil$\emph{a
weight} is a one dimensional representation of an algebra and are analogous to
group characters$\rfloor$ of the adjoint representation. The collection of
roots forms a \emph{root system} and much of the structure of $\mathcal{G}$
derives from its root system. Indeed, the \emph{classification} of complex
semisimple Lie algebras described above come from a classification of the
associated root systems, which in turn are classified by their Dynkin diagrams.

\begin{remark}
Cartan subalgebras and the associated root systems are basic tools for
understanding both the classification of semisimple Lie algebras and their
representation theory.
\end{remark}

\paragraph{Weyl group}

Let $\mathcal{G}$ be a semisimple Lie algebra, let $\mathcal{H}$ be a Cartan
subalgebra, and let $\mathcal{R}$ be the associated root system. For each
$\alpha\in\mathcal{R}$, we can consider the reflection $s_{\alpha}$ about the
hyperplane perpendicular to $\U{3b1} $. One of the basic properties of root
systems ensures that $\mathcal{R}$ is invariant under each $s_{\alpha}$. The
\emph{Weyl group} is then the group of linear transformations of $\mathcal{H}$
generated by the $s_{\alpha}$. The Weyl group is an important symmetry of the
problem; for example, the weights of any finite-dimensional representation of
$\mathcal{G}$ are invariant under the Weyl group.

\begin{example}
The case of $sl\left(  n,\mathbb{C}\right)  $
\end{example}

If $\mathcal{G}=sl\left(  n,\mathbb{C}\right)  $ then $\mathcal{H}$ may be
taken to be the diagonal subalgebra of $\mathcal{G}$ consisting of diagonal
matrices whose diagonal entries sum to zero. Since $\mathcal{H}$ has dimension
$n-1$, we see that $sl\left(  n,\mathbb{C}\right)  $ has rank $n-1$.

The root vectors $X$ in this case may be taken to be the matrices $E_{ij}$
with $i\neq j$, where $E_{ij}$ is the matrix with a $1$ in the $(i,j)$
position and zeros elsewhere. If $h$ is a diagonal matrix with diagonal
entries $\left\{  \U{3bb} _{1},\ldots\ ,\U{3bb} _{n}\right\}  $, then we have%

\begin{equation}
\left[  h,E_{ij}\right]  =\left(  \U{3bb} _{i}-\U{3bb} _{j}\right)  E_{ij}.
\end{equation}
Thus, the roots fors $sl\left(  n,\mathbb{C}\right)  $ are the linear
functionals
\begin{equation}
\alpha_{ij}\left(  h\right)  =\U{3bb} _{i}-\U{3bb} _{j}.
\end{equation}
After identifying $\mathcal{H}$ with its dual $\mathcal{H}^{d}$, the roots
become the vectors%
\begin{equation}
\alpha_{ij}=e_{i}-e_{j}%
\end{equation}
in the space of $n$-tuples that sum to zero. This is the root system known as
$A_{n-1}$ in the conventional labeling.\qquad

The reflection associated to the root $\alpha_{ij}$ acts on $\mathcal{H}$ by
transposing $i$ and $j$ diagonal entries. The Weyl group is then just the
permutation group on $n$ elements, acting by permuting the diagonal entries of
matrices in $\mathcal{H}$.

\paragraph{Role in the classification}

\emph{Semisimple Lie} algebras are ultimately classified by their \emph{root
systems} (which in turn are classified by their \emph{Dynkin diagrams}). The
argument is as follows.

\begin{itemize}
\item First, one proves that the \emph{Cartan subalgebra} of a semisimple Lie
algebra is unique up to isomorphism. It follows that the root system \emph{is
independent} (up to isomorphism) of the choice of Cartan subalgebra.

\item Next, one shows that the \emph{root system} determines the Lie algebra
up to isomorphism. (Thus, for example, there cannot be two nonisomorphic Lie
algebras both having the root system $A_{n}).$

\item Finally, one shows that every root system comes from a Lie algebra. This
can be done in a case-by-case fashion, with the only hard part being the
exceptional root systems of type $E,F,$ and $G$. Alternatively, one can use a
systematic procedure, using \emph{Serre's} relations.
\end{itemize}

\paragraph{Significance}

The significance of semisimplicity comes firstly from the \emph{Levi
decomposition}, which states that \emph{every finite dimensional Lie algebra}
is the \emph{semidirect product of a solvable ideal (its radical) and a
semisimple algebra}. In particular, there is \emph{no nonzero} Lie algebra
\emph{that is both solvable and semisimple}.

\begin{remark}
Semisimple Lie algebras have a very elegant classification, in stark contrast
to solvable Lie algebras. Semisimple Lie algebras over an algebraically closed
field are completely classified by their root system, which are in turn
classified by Dynkin diagrams. Semisimple algebras over non-algebraically
closed fields can be understood in terms of those over the algebraic closure,
though the classification is somewhat more intricate; see real form for the
case of real semisimple Lie algebras, which were classified by \'{E}lie Cartan.
\end{remark}

\begin{remark}
Further, the representation theory of semisimple Lie algebras is much cleaner
than that for general Lie algebras. For example, the Jordan decomposition in a
semisimple Lie algebra coincides with the Jordan decomposition in its
representation; this is not the case for Lie algebras in general.

\begin{remark}
If $\mathcal{G}$ is semisimple, then
\begin{equation}
\mathcal{G}=\left[  \mathcal{G},\mathcal{G}\right]  .
\end{equation}

\end{remark}

\begin{remark}
In particular, every \emph{linear semisimple Lie algebra} is a subalgebra of
$sl\left(  n,\mathbb{C}\right)  $ the special linear Lie algebra. The study of
the structure of $sl\left(  n,\mathbb{C}\right)  $ constitutes an important
part of the representation theory for semisimple Lie algebras.
\end{remark}
\end{remark}

\paragraph{Generalizations}

Semisimple Lie algebras admit certain generalizations. Firstly, many
statements that are true for \emph{semisimple Lie algebras} are true more
generally for \emph{reductive Lie algebras}.

Abstractly, \emph{a reductive Lie algebra is one whose adjoint representation
is completely reducible,} while concretely, \emph{a reductive Lie algebra is a
direct sum of a semisimple Lie algebra and an abelian Lie algebra}; for
example, $sl\left(  n,\mathbb{C}\right)  $ is semisimple, and $gl\left(
n,\mathbb{C}\right)  $ is reductive. Many properties of semisimple Lie
algebras depend only on reducibility.

\begin{remark}
Many properties of complex semisimple/reductive Lie algebras are true not only
for semisimple/reductive Lie algebras over algebraically closed fields, but
more generally for \emph{split semisimple/reductive} Lie algebras over other fields
\end{remark}

\begin{remark}
Semisimple/reductive Lie algebras over algebraically closed fields are always
split, but over other fields this is not always the case. Split Lie algebras
have essentially the same representation theory as semsimple Lie algebras over
algebraically closed fields, for instance, the splitting Cartan subalgebra
playing the same role as the Cartan subalgebra plays over algebraically closed
fields. This is the approach followed in (Bourbaki 2005), for instance, which
classifies representations of split semisimple/reductive Lie algebras.
\end{remark}


\end{document}